% Mathematical TeX extracted from the cited web pages, reformatted by the assistant. Not the native full chapter.

\begin{lemma}[30.18.1, Chow's Lemma, Tag 0200]
Let $S$ be a Noetherian scheme. Let $f:X\to S$ be a separated morphism of finite type. Then there exist an $n\geq0$ and a diagram
\[
\xymatrix{X\ar[rd]&X'\ar[d]\ar[l]^\pi\ar[r]&\mathbf P^n_S\ar[dl]\\&S&}
\]
where $X'\to\mathbf P^n_S$ is an immersion, and $\pi:X'\to X$ is proper and surjective. Moreover, we may arrange it such that there exists a dense open subscheme $U\subset X$ such that $\pi^{-1}(U)\to U$ is an isomorphism.
\end{lemma}
\begin{proof}
All of the schemes we will encounter during the rest of the proof are going to be of finite type over the Noetherian scheme $S$ and hence Noetherian (see Morphisms, Lemma 29.15.6). All morphisms between them will automatically be quasi-compact, locally of finite type and quasi-separated, see Morphisms, Lemma 29.15.8 and Properties, Lemmas 28.5.4 and 28.5.8.

The scheme $X$ has only finitely many irreducible components (Properties, Lemma 28.5.7). Say $X=X_1\cup\ldots\cup X_r$ is the decomposition of $X$ into irreducible components. Let $\eta_i\in X_i$ be the generic point. For every point $x\in X$ there exists an affine open $U_x\subset X$ which contains $x$ and each of the generic points $\eta_i$. See Properties, Lemma 28.30.4. Since $X$ is quasi-compact, we can find a finite affine open covering $X=U_1\cup\ldots\cup U_m$ such that each $U_i$ contains $\eta_1,\ldots,\eta_r$. In particular we conclude that the open $U=U_1\cap\ldots\cap U_m\subset X$ is a dense open. This and the fact that the $U_i$ are affine opens covering $X$ are all that we will use below.

Let $X^*\subset X$ be the scheme theoretic closure of $U\to X$, see Morphisms, Definition 29.6.2. Let $U_i^*=X^*\cap U_i$. Note that $U_i^*$ is a closed subscheme of $U_i$. Hence $U_i^*$ is affine. Since $U$ is dense in $X$ the morphism $X^*\to X$ is a surjective closed immersion. It is an isomorphism over $U$. Hence we may replace $X$ by $X^*$ and $U_i$ by $U_i^*$ and assume that $U$ is scheme theoretically dense in $X$, see Morphisms, Definition 29.7.1.

By Morphisms, Lemma 29.40.3 we can find an immersion $j_i:U_i\to\mathbf P_S^{n_i}$ for each $i$. By Morphisms, Lemma 29.7.7 we can find closed subschemes $Z_i\subset\mathbf P_S^{n_i}$ such that $j_i:U_i\to Z_i$ is a scheme theoretically dense open immersion. Note that $Z_i\to S$ is proper, see Morphisms, Lemma 29.44.5. Consider the morphism
\[
j=(j_1|_U,\ldots,j_m|_U):U\longrightarrow
\mathbf P_S^{n_1}\times_S\ldots\times_S\mathbf P_S^{n_m}.
\]
By the lemma cited above we can find a closed subscheme $Z$ of $\mathbf P_S^{n_1}\times_S\ldots\times_S\mathbf P_S^{n_m}$ such that $j:U\to Z$ is an open immersion and such that $U$ is scheme theoretically dense in $Z$. The morphism $Z\to S$ is proper. Consider the $i$th projection
\[\operatorname{pr}_i|_Z:Z\longrightarrow\mathbf P_S^{n_i}.\]
This morphism factors through $Z_i$ (see Morphisms, Lemma 29.6.6). Denote $p_i:Z\to Z_i$ the induced morphism. This is a proper morphism, see Morphisms, Lemma 29.42.7 for example. At this point we have that $U\subset U_i\subset Z_i$ are scheme theoretically dense open immersions. Moreover, we can think of $Z$ as the scheme theoretic image of the ``diagonal'' morphism $U\to Z_1\times_S\ldots\times_S Z_m$.

Set $V_i=p_i^{-1}(U_i)$. Note that $p_i|_{V_i}:V_i\to U_i$ is proper. Set $X'=V_1\cup\ldots\cup V_m$. By construction $X'$ has an immersion into the scheme $\mathbf P_S^{n_1}\times_S\ldots\times_S\mathbf P_S^{n_m}$. Thus by the Segre embedding (see Constructions, Lemma 27.13.6) we see that $X'$ has an immersion into a projective space over $S$.

We claim that the morphisms $p_i|_{V_i}:V_i\to U_i$ glue to a morphism $X'\to X$. Namely, it is clear that $p_i|_U$ is the identity map from $U$ to $U$. Since $U\subset X'$ is scheme theoretically dense by construction, it is also scheme theoretically dense in the open subscheme $V_i\cap V_j$. Thus we see that $p_i|_{V_i\cap V_j}=p_j|_{V_i\cap V_j}$ as morphisms into the separated $S$-scheme $X$, see Morphisms, Lemma 29.7.10. We denote the resulting morphism $\pi:X'\to X$.

We claim that $\pi^{-1}(U_i)=V_i$. Since $\pi|_{V_i}=p_i|_{V_i}$ it follows that $V_i\subset\pi^{-1}(U_i)$. Consider the diagram
\[
\xymatrix{V_i\ar[r]\ar[rd]_{p_i|_{V_i}}&\pi^{-1}(U_i)\ar[d]\\&U_i}
\]
Since $V_i\to U_i$ is proper we see that the image of the horizontal arrow is closed, see Morphisms, Lemma 29.42.7. Since $V_i\subset\pi^{-1}(U_i)$ is scheme theoretically dense (as it contains $U$) we conclude that $V_i=\pi^{-1}(U_i)$ as claimed.

This shows that $\pi^{-1}(U_i)\to U_i$ is identified with the proper morphism $p_i|_{V_i}:V_i\to U_i$. Hence we see that $X$ has a finite affine covering $X=\bigcup U_i$ such that the restriction of $\pi$ is proper on each member of the covering. Thus by Morphisms, Lemma 29.42.3 we see that $\pi$ is proper.

Finally we have to show that $\pi^{-1}(U)=U$. To see this we argue in the same way as above using the diagram
\[\xymatrix{U\ar[r]\ar[rd]&\pi^{-1}(U)\ar[d]\\&U}\]
and using that $\id_U:U\to U$ is proper and that $U$ is scheme theoretically dense in $\pi^{-1}(U)$.
\end{proof}

